\section{Local removal of interaction energy}
\label{sec:energy}

The estimate on zero modes controls the cost of a missing electron.  We now
connect an arbitrary state to the zero space by a local process that removes
particles.  Its interaction energy decays exponentially, and its expected
particle loss is bounded by its initial interaction energy.  In particular,
adding electrons above Laughlin filling has a uniform energy cost.

Write $B_p=B(v_p)$, $0\le p\le 2q-2$, for the normalized pair annihilators
of \eqref{eq:pair-rows}, and put $d_a=2q-1$.  For a rotation $g\in SU(2)$,
$A(g)$ denotes the conjugate of an operator $A$ by the Fock representation.
Haar measure $dg$ has mass one.  The rotation average of any normalized
pair row resolves the interaction:
\begin{equation}\label{eq:pair-resolution}
 d_a\int B_p(g)^*B_p(g)\,dg=H_q.
\end{equation}
This is Schur's lemma on the spin-$(q-1)$ pair space $V_q$.

\begin{theorem}[Local energy descent]\label{thm:local-energy-descent}
There are constants $c>0$, $L<\infty$, and $q_0$ such that, for every
$q\ge q_0$, there is a measurable family of rapidly local physical
operators $J_y$ on $\mathcal F_q$ satisfying
\begin{equation}\label{eq:energy-descent}
 \int J_y^*J_y\,d\nu_q(y)=H_q,
 \qquad
 \int J_y^*H_qJ_y\,d\nu_q(y)\le H_q^2-cH_q.
\end{equation}
Every $J_y$ has a definite particle loss $r_y\in\{2,\ldots,L\}$, so that
$[\mathcal N,J_y]=-r_yJ_y$.  The center measure has bounded mass per unit
ball, and all rapid-locality bounds are independent of $q$ and of the
particle sector.
\end{theorem}

The first operation in the construction is $B_0(g)$, which removes a pair
near the point $g\cdot\mathrm{north}$.  We then partially empty a fixed
neighborhood of that point.  This second operation suppresses the nearby
pair correlations left behind by the first removal.  We first identify the
part of $H_q^2$ that pays for the surviving correlations.

\subsection{The four-body term available for descent}

The maps $\mathcal L_k$ and $W_4$ are defined in
\eqref{eq:lift} and \eqref{eq:wedge-maps}.  Thus
$\mathcal L_4(|u\rangle\langle v|)=B(u)^*B(v)$ and
$W_4(v\otimes w)=v\wedge w$.
Let $P_{q,\ell}$ be the projection in $V_q\otimes V_q$ onto total spin
$2(q-1)-\ell$, for $0\le\ell\le2q-2$, and set
\begin{equation}\label{eq:high-pair-blocks}
 R_q^{\mathrm{hi}}=\sum_{\ell=23}^{2q-2}P_{q,\ell},
 \qquad R_q^{\mathrm{lo}}=I-R_q^{\mathrm{hi}}.
\end{equation}
For large $q$, $R_q^{\mathrm{lo}}$ contains precisely the small deficits
$\ell=0,\ldots,22$.
Proposition~\ref{prop:retained-blocks}, proved in the appendix, gives
\begin{equation}\label{eq:retained-energy}
 H_q^2\ge c_0H_q+
 \mathcal L_4(W_4R_q^{\mathrm{hi}}W_4^*)
 \qquad(0<c_0<1/25)
\end{equation}
for all sufficiently large $q$, uniformly on $\mathcal F_q$.

This statement contains more information than the unperturbed spectral
gap.  The normal-ordering identity for $H_q^2$ contains the full four-body
term $\mathcal L_4(W_4W_4^*)$.  The finite comparison of
\cite[Sections 2--6]{Gap} uses only the blocks $\ell\le22$.
The appendix includes the finite certificate and its transfer to finite
$q$, and keeps the other blocks with coefficient one.  The loss construction will use precisely this
part of the square.

\subsection{Emptying a neighborhood without increasing distant pair energy}

Choose a smooth nondecreasing $T:\mathbb R\to[0,1]$ such that $T=0$ on
$(-\infty,1]$ and $T=1$ on $[3,\infty)$.  We choose it flat at the two
endpoints and, in addition, require
\begin{equation}\label{eq:evacuation-cutoff}
 |T^{(r)}(x)|\le C_r\sqrt{1-T(x)}\qquad(r\ge1).
\end{equation}
For example, with $\eta(x)=e^{-1/x}$ for $x>0$ and $\eta(x)=0$ for $x\le0$,
one may take
$T(x)=\eta(x-1)/(\eta(x-1)+\eta(3-x))$.
The estimate follows by differentiating near $x=3$; an exponential
$e^{-1/(3-x)}$ times any fixed negative power of $3-x$ is bounded by a
constant times $e^{-1/(2(3-x))}$.  Away from that endpoint the estimate is
immediate from smoothness.

For an integer $M\ge2$, put $t_j=T(j/M)$ and $n_j=c_j^*c_j$.  The two
operators
\begin{equation}\label{eq:single-mode-loss}
 K_{j,0}=I-n_j+t_jn_j,
 \qquad K_{j,1}=\sqrt{1-t_j^2}\,c_j
\end{equation}
define a trace-preserving channel, since
$K_{j,0}^*K_{j,0}+K_{j,1}^*K_{j,1}=I$.
Compose these channels in increasing order of $j$, for $0\le j<3M$.
Write $K_\alpha$ for the resulting products of Kraus operators, where
$\alpha\in\{0,1\}^{3M}$, and write $\mathcal E_M$ for the channel.
Its dual on observables is
$\mathcal E_M^*(O)=\sum_\alpha K_\alpha^*OK_\alpha$.

The canonical anticommutation relations imply a useful exact formula.
On an even normal-ordered monomial, the dual of the $j$th channel
multiplies by $t_j$ for each occurrence of $c_j$ or $c_j^*$.
Indeed, if neither occurs, the even monomial commutes with both operators
of mode $j$ and completeness gives the claim.  If exactly one occurs,
the loss contribution vanishes and the no-loss contribution gives $t_j$.
If both occur, writing them as an occupation factor gives $t_j^2$.
Consequently
\begin{equation}\label{eq:evacuated-interaction}
 \mathcal E_M^*(H_q)=\sum_{b=1}^{2q-1}\widetilde B_{b-1}^*
                                      \widetilde B_{b-1},
 \qquad
 \widetilde B_{b-1}
   =\sum_{\substack{i<j\\i+j=b}}t_it_j\,
                   \overline{v_{b-1}(i,j)}\,c_jc_i.
\end{equation}
Here and below terms with an orbital index outside $\{0,\ldots,q\}$ are
absent.

\begin{lemma}[Pair energy after local loss]\label{lem:evacuation-rows}
There is a sequence $\varepsilon_M\to0$ such that, for every $M\ge2$
and every $q\ge30M$,
\begin{equation}\label{eq:evacuation-rows}
 \sum_{b=1}^{2q-1}\widetilde B_{b-1}^*\widetilde B_{b-1}
 \le\sum_{b>M}B_{b-1}^*B_{b-1}+\varepsilon_MI
 \quad\hbox{on }\mathcal F_q.
\end{equation}
\end{lemma}
\begin{proof}
The following two elementary estimates specialize the general pair
localization bounds to the cutoff used here.  In particular, the
exponential estimate on distant rows will let us preserve coefficient
one on their energy.  Both follow directly from the coefficients in
\eqref{eq:pair-rows}.  For $b\le q/2$ and every fixed integer $r\ge0$,
\begin{equation}\label{eq:pair-absolute-moments}
 \sum_{\substack{i<j\\i+j=b}}|v_{b-1}(i,j)|\,|i-b/2|^r
 \le C_r b^{r/2+1/4}.
\end{equation}
For $b>10M$ and $q\ge30M$,
\begin{equation}\label{eq:pair-remote-tail}
 \sum_{\substack{i<j,\ i+j=b\\i<3M}}|v_{b-1}(i,j)|
 \le Cb^{3/2}e^{-c_1b}.
\end{equation}
Constants may change after decreasing $c_1>0$.
Here is a direct verification.  Set
\[
 p_{q,b}(i)=\frac{\binom qi\binom q{b-i}}{\binom{2q}b}.
\]
The exact pair formula gives
\begin{equation}\label{eq:pair-hypergeometric}
 |v_{b-1}(i,b-i)|^2
 =\frac{2(2q-1)(2i-b)^2}{b(2q-b)}p_{q,b}(i),
 \qquad
 \frac{p_{q,b}(i+1)}{p_{q,b}(i)}
 =\frac{(q-i)(b-i)}{(i+1)(q-b+i+1)}.
\end{equation}
When $b\le q/2$, these ratios, their inverses, and symmetry about $b/2$
show
\[
 p_{q,b}(i)\le Cb^{-1/2}
                  \exp\{-c_2(i-b/2)^2/b\}.
\]
For completeness, on the upper half $i\ge(b-1)/2$, the first factor
in the ratio is at most one and the second is at most
$\exp\{-(2i+1-b)/(b+1)\}$.  Summing logarithms from a middle
index gives the Gaussian upper bound relative to the middle value.  In the range $|i-b/2|\le\sqrt b$, the absolute value of
that logarithm is at most $C(1+|i-b/2|)/b$.  Summing these lower bounds and
using $\sum_i p_{q,b}(i)=1$ bounds the middle value by $C/\sqrt b$.
Insertion in \eqref{eq:pair-hypergeometric}, followed by summing Gaussian
moments, proves \eqref{eq:pair-absolute-moments}.

To prove \eqref{eq:pair-remote-tail}, there is nothing to check if
$b>q+3M$, because the sum is empty.  Otherwise $b\le q+3M\le11q/10$,
so $2(2q-1)/[b(2q-b)]\le C/b$ in
\eqref{eq:pair-hypergeometric}.
For every allowed $i\le2b/5$ and $b\ge20$, the ratio in
\eqref{eq:pair-hypergeometric} is at least $4/3$.
Starting with $i<3M<3b/10$ and taking at least $b/10-O(1)$ steps toward
$2b/5$ gives $p_{q,b}(i)\le Ce^{-c b}$, since each probability is at most
one.  The square root of \eqref{eq:pair-hypergeometric} is therefore at
most $C\sqrt b\,e^{-c_1b}$, and there are at most $b$ terms.  This proves
\eqref{eq:pair-remote-tail}.

The modified row vanishes when $b\le M$.  For $M<b\le10M$, put
$x=b/(2M)$ and $t=T(x)$.  The smooth even function
$z\mapsto T(x+z)T(x-z)$ has zero first derivative at zero.  By
\eqref{eq:evacuation-cutoff}, every derivative of positive order there is
bounded by $C_r\sqrt{1-t}$.  Taylor expansion through degree $41$, with a
bounded remainder of order $42$, and
\eqref{eq:pair-absolute-moments} yield
\begin{equation}\label{eq:weighted-row-error}
 \widetilde B_{b-1}=t^2 B_{b-1}+E_b,
 \qquad
 \|E_b\|\le CM^{-3/4}\sqrt{1-t}+CM^{-20}.
\end{equation}
We used $\|c_jc_i\|\le1$.  More explicitly, the degree-$r$ contribution
is bounded by
\[
 C_rM^{-r}b^{r/2+1/4}\sqrt{1-t},\qquad 2\le r\le41,
\]
and the remainder is at most $CM^{-42}b^{21+1/4}$.

Let $\delta=M^{-8}$ and $s=1-t+\delta$.  For any vector $\phi$, Cauchy's
inequality gives
\[
 \|(t^2B_{b-1}+E_b)\phi\|^2
 \le(t^4+s)\|B_{b-1}\phi\|^2
       +(1+t^4/s)\|E_b\phi\|^2.
\]
Since $t^4+s\le1+\delta$, \eqref{eq:weighted-row-error} bounds the second
term by $CM^{-3/2}\|\phi\|^2$.  Summing over the $O(M)$ central rows
costs $CM^{-1/2}\|\phi\|^2$.
The additional coefficient $\delta$ also disappears in this error:
\eqref{eq:pair-absolute-moments} with $r=0$ implies
$\sum_{M<b\le10M}\|B_{b-1}\|^2\le CM^{3/2}$.

Finally, for $b>10M$ the coefficients of
$\widetilde B_{b-1}-B_{b-1}$ are supported where $i<3M$ or $j<3M$.
The increasing ordering reduces this to $i<3M$.
Equation~\eqref{eq:pair-remote-tail} bounds this difference in operator
norm by $Cb^{3/2}e^{-c_1b}$, while normalization of the pair vector gives
$\|B_{b-1}\|\le\sqrt b$.  The norm of the difference of the two row
squares is therefore summable over $b>10M$, with sum at most
$Ce^{-c_3M}$.  Combining the three ranges proves the assertion, with
$\varepsilon_M\le C(M^{-1/2}+M^{-13/2}+e^{-c_3M})$.
\end{proof}

\subsection{Rotation averaging suppresses the small-deficit blocks}

Retaining only the rows whose index sum exceeds $M$ gives a positive
operator on the two pair factors:
\begin{equation}\label{eq:pair-tail-average}
 A_{q,M}=d_a\int
 |v_0(g)\rangle\langle v_0(g)|\otimes
 \left(\sum_{k+1>M}|v_k(g)\rangle\langle v_k(g)|\right)\,dg.
\end{equation}
If the sum includes all $k$, the operator equals $I$ by Schur's lemma.
Consequently $0\le A_{q,M}\le I$.  This bound handles every retained
block with the same coefficient one as in \eqref{eq:retained-energy}.
It remains to make the coefficients on the finitely many unretained
blocks small and bound their lifts by $H_q$.

\begin{lemma}[Vanishing weight on fixed deficits]\label{lem:pair-spin-tail}
For $A_{q,M}$ in \eqref{eq:pair-tail-average}, there are coefficients
$0\le a_{q,M,\ell}\le1$ such that
\[
 A_{q,M}=\sum_{\ell=0}^{2q-2}a_{q,M,\ell}P_{q,\ell}.
\]
For every fixed $\ell\ge0$,
\begin{equation}\label{eq:spin-tail-limit}
 \lim_{M\to\infty}\limsup_{q\to\infty}a_{q,M,\ell}=0.
\end{equation}
Moreover, for each fixed $\ell$ there is $C_\ell<\infty$ such that
\begin{equation}\label{eq:fixed-block-energy}
 \mathcal L_4(W_4P_{q,\ell}W_4^*)\le C_\ell H_q
\end{equation}
for all sufficiently large $q$.
\end{lemma}
\begin{proof}
The tensor product of two spin-$(q-1)$ representations is multiplicity
free.  Hence rotation invariance and $0\le A_{q,M}\le I$ prove the first
assertion.  Put $a=q-1$ and $d_{q,\ell}=4q-3-2\ell$.
If $C_{q}(\ell,k)$ is the coupling coefficient of $v_0\otimes v_k$ in
total spin $2a-\ell$ at total lowering degree $k$, then
\begin{equation}\label{eq:spin-tail-coefficient}
 a_{q,M,\ell}=\frac{2q-1}{d_{q,\ell}}
                    \sum_{k+1>M}|C_q(\ell,k)|^2.
\end{equation}
The unrestricted sum, including all $k$, has coefficient exactly one.

We give the fixed-$k$ limit, including its normalization.  At a fixed
total degree, the two spin lowering operators, divided by $\sqrt{2a}$,
converge to multiplication by two independent oscillator variables
$X,Y$.  The highest vector of deficit $\ell$ converges to
$(X-Y)^\ell/\sqrt{2^\ell\ell!}$: applying the raising operator gives
successive coefficient ratio
$-\sqrt{(\ell-p)/(p+1)}$ in the limit, which fixes this normalized vector.
Normalized subsequent lowering multiplies by
$(X+Y)/\sqrt{2(k-\ell)}$ at degree $k$.  Thus the degree-$k$ coupled
vector converges, in the orthonormal monomial basis, to
\[
 \frac{(X-Y)^\ell(X+Y)^{k-\ell}}
      {2^{k/2}\sqrt{\ell!(k-\ell)!}}.
\]
Its coefficient of $Y^k/\sqrt{k!}$ has squared modulus
$2^{-k}\binom{k}{\ell}$, and hence
\[
 |C_q(\ell,k)|^2\longrightarrow2^{-k}\binom{k}{\ell}.
\]
There are only finitely many degrees below a fixed $M$.  Subtracting
these from the exact unrestricted coefficient one gives
\[
 \lim_{q\to\infty}a_{q,M,\ell}
 =1-\frac12\sum_{\ell\le k\le M-1}2^{-k}\binom{k}{\ell}.
\]
The identity
$\sum_{k=\ell}^{\infty}2^{-k}\binom{k}{\ell}=2$
proves \eqref{eq:spin-tail-limit}.  This finite-head argument requires no
uniform convergence over an unbounded set of coupling coefficients.

For \eqref{eq:fixed-block-energy}, choose a normalized highest vector
$z_{q,\ell}=\sum_{p=0}^{\ell}s_pv_p\otimes v_{\ell-p}$, with
$\sum_p|s_p|^2=1$.  Its wedge annihilator is
$\sum_p\overline{s_p}B_{\ell-p}B_p$.
Every $B_{\ell-p}$ uses only the first $\ell+2$ orbitals and has norm at
most $\sqrt{\ell+2}$.  Therefore
\[
 \|B(W_4z_{q,\ell})\psi\|^2
 \le(\ell+2)\sum_{p=0}^{\ell}\|B_p\psi\|^2.
\]
Resolve $P_{q,\ell}$ by the orbit of its highest vector and use
\eqref{eq:pair-resolution}.  This yields
\[
 \mathcal L_4(W_4P_{q,\ell}W_4^*)
 \le\frac{d_{q,\ell}}{2q-1}(\ell+2)(\ell+1)H_q,
\]
whose coefficient is bounded for fixed $\ell$.
\end{proof}

\begin{proof}[Proof of Theorem~\ref{thm:local-energy-descent}]
For the moment fix $M$, and define
\[
 J_{g,\alpha}=K_\alpha(g)B_0(g),
 \qquad d\nu_q(g,\alpha)=(2q-1)\,dg
\]
with counting measure in $\alpha$.
Completeness of the loss channel and \eqref{eq:pair-resolution} give
$\int J_y^*J_y\,d\nu_q=H_q$.
Lemma~\ref{lem:evacuation-rows}, applied to $B_0(g)\psi$ and then averaged,
gives
\begin{align}
 \int J_y^*H_qJ_y\,d\nu_q(y)
 &=d_a\int B_0(g)^*\mathcal E_{M,g}^*(H_q)B_0(g)\,dg\notag\\
 &\le\mathcal L_4(W_4A_{q,M}W_4^*)+\varepsilon_MH_q.
 \label{eq:loss-four-body}
\end{align}
The equality between the averaged pair products and the lifted four-body
operator uses ordered tensor factors.  There is no factor of two:
$B(v\wedge w)=B(w)B(v)$ and the lift sums each ordered pair once.

By Lemma~\ref{lem:pair-spin-tail} and positivity of the lift,
\[
 \mathcal L_4(W_4A_{q,M}W_4^*)
 \le\mathcal L_4(W_4R_q^{\mathrm{hi}}W_4^*)
       +\left(\sum_{\ell=0}^{22}C_\ell a_{q,M,\ell}\right)H_q.
\]
Fix $c_0=1/50$ in \eqref{eq:retained-energy}.  Choose $M$ large enough,
and then $q$ large enough, that
$\varepsilon_M+\sum_{\ell=0}^{22}C_\ell a_{q,M,\ell}\le c_0/2$.
Equations~\eqref{eq:retained-energy} and \eqref{eq:loss-four-body} prove
\eqref{eq:energy-descent} with $c=c_0/2$.

Each $K_\alpha$ is a product of $3M$ operators of particle loss zero or
one, so $J_{g,\alpha}$ has loss between $2$ and $L=3M+2$.
All modes involved in $J_{g,\alpha}$ have index below $3M$, apart from
identity factors.  To see operator-norm locality directly, use the frame
isometry $\mathsf S_q:U_q\to\ell^2(\Lambda_q)$ and write $c(f)=B(f)$ on the slot
Fock space.  If $P_R$ cuts off slots outside the radius-$R$ ball around
$g\cdot\mathrm{north}$, the canonical anticommutation relations give
\[
 \|c(\mathsf S_qe_j(g))-c(P_R\mathsf S_qe_j(g))\|
   =\|(I-P_R)\mathsf S_qe_j(g)\|_2\le C_{D,M}R^{-D}.
\]
The last estimate follows from Lemma~\ref{locality:frame} and the slot
volume bound, with a larger decay order before summing squares.
The same estimate for a product follows by telescoping its finitely many
factors, or by Lemma~\ref{locality:calculus}.  Thus the jumps are rapidly
local around $g\cdot\mathrm{north}$, uniformly in $q$.
There are only $2^{3M}$ auxiliary labels, with $M$ now fixed.
The pushforward of $(2q-1)dg$ to the physical sphere has area density
$(2q-1)/(2\pi q)$; the remaining rotation angle has mass one.
Including all auxiliary labels therefore gives total center density
$2^{3M}(2q-1)/(2\pi q)$, uniformly bounded in $q$ for this fixed $M$.
These observations prove all locality and density assertions.
\end{proof}

\subsection{Particle loss, excess charge, and a coercive Fock Hamiltonian}

Let $q=3(N-1)$ and let $\rho_t$ evolve by the trace-preserving equation
whose dual generator is
\begin{equation}\label{eq:loss-generator}
 \mathcal D_q(O)=\int J_y^*OJ_y\,d\nu_q(y)-\tfrac12\{H_q,O\}.
\end{equation}
This is a finite-dimensional Lindblad generator~\cite{Lindblad}.  Positivity and trace
preservation also follow directly by iterating variation of constants,
with the completely positive gain map
$\rho\mapsto\int J_y\rho J_y^*\,d\nu_q(y)$ and the no-jump map
$\rho\mapsto e^{-tH_q/2}\rho e^{-tH_q/2}$.
We write $\langle O\rangle_t=\operatorname{Tr}(\rho_tO)$.

\begin{proposition}[Control of particle number]\label{prop:grand-canonical-coercivity}
There are constants $C_{\mathrm{num}}<\infty$, $\mu>0$, $a>0$, and
$C_G<\infty$, independent of sufficiently large $q=3(N-1)$, such that
\begin{equation}\label{eq:excess-number-bound}
 \mathcal N-N\le C_{\mathrm{num}}H_q.
\end{equation}
The operator
\begin{equation}\label{eq:coercive-G}
 G_q=H_q+\mu(N-\mathcal N)
\end{equation}
satisfies
\begin{equation}\label{eq:G-coercivity}
 G_q\ge a\bigl(H_q+|N-\mathcal N|\bigr)
 \ge g\bigl(I-P_{\mathrm L,N}\bigr)
\end{equation}
for a constant $g>0$, with $P_{\mathrm L,N}$ extended by zero on the
other Fock sectors.  For every initial density matrix, the evolution
\eqref{eq:loss-generator} obeys
\begin{align}
 \langle H_q\rangle_t&\le e^{-ct}\langle H_q\rangle_0,
       \label{eq:loss-energy-decay}\\
 0\le\langle\mathcal N\rangle_0-\langle\mathcal N\rangle_t
    &\le C_{\mathrm{num}}\langle H_q\rangle_0,
       \label{eq:loss-number-control}\\
 \langle G_q\rangle_t&\le C_G\langle G_q\rangle_0.
       \label{eq:loss-G-control}
\end{align}
Every limit point $\rho_\infty$ is supported on $\ker H_q$, and
\begin{equation}\label{eq:loss-hole-control}
 0\le\operatorname{Tr}\rho_\infty(N-\mathcal N)
       \le C_G\mu^{-1}\langle G_q\rangle_0.
\end{equation}
\end{proposition}
\begin{proof}
Theorem~\ref{thm:local-energy-descent} gives
$\mathcal D_q(H_q)\le-cH_q$, proving
\eqref{eq:loss-energy-decay}.  Since $[\mathcal N,J_y]=-r_yJ_y$,
\[
 -\mathcal D_q(\mathcal N)=\int r_yJ_y^*J_y\,d\nu_q(y),
 \qquad 2H_q\le-\mathcal D_q(\mathcal N)\le LH_q.
\]
Integration proves \eqref{eq:loss-number-control} with
$C_{\mathrm{num}}=L/c$.
Compactness of the density matrices in each fixed finite-dimensional
Fock space gives limit points as $t\to\infty$.  By energy decay each is
supported on $\ker H_q$.  Lemma~\ref{lem:zero-modes} shows that this kernel
contains no sector of particle number exceeding $N$.  Applying
\eqref{eq:loss-number-control} to a state of definite particle number $n$
and passing to a limit point gives $n-N\le C_{\mathrm{num}}\langle H_q\rangle_0$.
This proves \eqref{eq:excess-number-bound} on each sector, and hence on
Fock space.

Choose $0<\mu\le\min(1,(2C_{\mathrm{num}})^{-1})$.
In sectors $n\le N$, $G_q=H_q+\mu(N-n)$ dominates
$\mu(H_q+N-n)$.  In sectors $n>N$, \eqref{eq:excess-number-bound} gives
$G_q\ge H_q/2$ and $H_q+n-N\le(1+C_{\mathrm{num}})H_q$.
This proves the first inequality in \eqref{eq:G-coercivity} with
$a=\min(\mu,[2(1+C_{\mathrm{num}})]^{-1})$.
For $n\ne N$, $|N-n|\ge1$; for $n=N$, the uniform gap and uniqueness in
Theorem~\ref{thm:unperturbed} give
$H_{N,q}\ge(1/25)(I-P_{\mathrm L,N})$.
Thus the second inequality holds with $g=a/25$.

Finally, energy decay and number control give
\[
 \langle G_q\rangle_t
 \le\langle G_q\rangle_0+\mu C_{\mathrm{num}}\langle H_q\rangle_0
 \le(1+\mu C_{\mathrm{num}}/a)\langle G_q\rangle_0.
\]
Take $C_G=1+\mu C_{\mathrm{num}}/a$.  Passing to a limit point, where
$H_q=0$, proves \eqref{eq:loss-hole-control}.
\end{proof}
